% Deep mean-field closure, 22 October 2026.
% Input after the bulk-tail, RKHS-certificate, and generic-atom fragments.
\section{The deep mean-field endpoint}
\label{sec:meanfield-generic-closure}
Throughout this section $\phi(0)=0$ and $\phi$ is $\lambda$-Lipschitz,
where $\lambda\ge1$ is fixed. The Bayesian model uses this same
activation at every neuron. The finite reconstructions used only for
counting may have different zero-at-zero $\lambda$-Lipschitz activations
at different neurons; the preceding capacity theorem covers precisely
this enlarged class. No claim that those auxiliary activations are
additional parameters of the Bayesian learner is made.

We now assemble the local covariance estimates into a finite network,
working backward to choose error tolerances and forward to reconstruct
the layers. At $\gamma=1$, let $H=L-1$, $m=2^d$, and write
\[
 Q_0(x,y)=d^{-1}\bar x\cdot\bar y,\qquad
 Q_s(x,y)=N^{-1}\sum_{i=1}^Nh_i^s(x)h_i^s(y).
\]
Conditionally on the preceding layers, the $N$ preactivation fields at
layer $s$ are independent Gaussians with covariance $Q_{s-1}$. This is
exact at every finite width.

\subsection{One-step reconstruction and its polynomial modulus}
Here a modulus specifies how accurately the incoming covariance must be
reconstructed to achieve a prescribed outgoing accuracy. Its reciprocal
will be polynomially bounded, allowing iteration at fixed depth without
an exponential dependence on $m$. Fix $K\ge1$, put
\[
 q=4K,\qquad A=4(q+\log2),
\]
and suppose the input covariance has diagonal at most $M\ge1$.
For a desired output accuracy $0<\alpha\le1$, set
\[
 B=\sqrt{2M\log(4m)},\qquad
 T=\max\{1,2B,96\lambda^2BMA/\alpha\}.
\]
Use the smooth cutoff from the preceding section. Mollify $\phi$ at
scale $\tau\le1$ small enough that the maximum-norm difference between
its cutoff covariance and that using the actual activation is at most
$\alpha/16$, uniformly in the Gaussian covariance. For example
$\tau=\min\{1,\alpha/[64C\lambda^2(T+1)]\}$ works with the absolute
constant $C$ in the smoothing estimate. Let $F_0,F_1,F_2$ be the bounded
observable and derivative budgets in that section, and put
\[
 G=16q/\alpha,\qquad
 s=\max\{1,\lceil G^2F_0^2\rceil\},
\]
\[
 C_G=\tfrac12\{F_2+2GF_1^2+2F_0(GF_2+G^2F_1^2)\},\qquad
 \eta=\frac{\alpha}{512A\lambda^2(\sqrt M+1)},
\]
\[
 r(\alpha)=\min\{1,\eta^2,\alpha/[16(1+C_G)]\},\qquad
 N_0(\alpha)=\left\lceil\frac{2s\log(2m^2+1)}q\right\rceil.
\]
Every displayed quantity is bounded by a polynomial in
$K,M,\log(2m),\alpha^{-1}$, or, for $r$, its reciprocal is so bounded.
The constants may depend on $\lambda$ and the fixed smooth cutoff.

For $N\ge N_0(\alpha)$, except with probability at most $2e^{-NK}$,
the empirical outgoing covariance is within $3\alpha/16$ of
\[
 \mathcal F_\Gamma(Q)+R,\qquad
 \|\Gamma\|_{1,\mathrm{ent}}\le G,
\]
where $R$ is a mixture of unit-RKHS auxiliary-activation atoms with total
coefficient mass at most $A$. Indeed, the tail decomposition error is at
most $\alpha/16$, smoothing costs at most $\alpha/16$, and the entropy
certificate costs at most $q/G=\alpha/16$. The projection failure is
at most $e^{-Nq}$; the entropy failure is at most
\[
 \exp\{s\log(2m^2+1)-N(q-1/2)\}\le e^{-NK}.
\]
These bounds hold conditionally and uniformly in the input covariance.

Now suppose $Q'$ is a positive semidefinite covariance with
$\|Q'-Q\|_{\max}\le r(\alpha)$, represented by finitely many weighted
features from the preceding layer. The selected $\Gamma$ and atom
certificates remain fixed while $Q$ is replaced by $Q'$. The bulk change
is at most $\alpha/16$. The generic atom construction, with target
accuracy $\alpha/8$, uses
\[
 n_A=\left\lceil8192A^2\lambda^4M^2\alpha^{-2}\log(4m^2)\right\rceil
\]
auxiliary neurons and approximates $R$ within $\alpha/8$ at covariance
$Q'$. Its certificate parameter is exactly the $\eta$ above.

To realize the bulk, sample independent Gaussian rows of covariance
$Q'$ tilted by the bounded smooth observable. Evaluate the actual
activation $\phi$ at those rows and absorb the scalar cutoff into their
covariance coefficients. Each retained cutoff feature is bounded by
$B_b=\lambda\sqrt5T$. Take
\[
 n_B=\left\lceil1024B_b^4\alpha^{-2}\log(16m^2)\right\rceil.
\]
Hoeffding's inequality and a union bound over $m^2$ entries give a
realization with empirical covariance within $\alpha/16$ of its own
expectation. The expectation differs from
$\mathcal F_\Gamma(Q')$ by at most $\alpha/16$ because of the smooth
activation approximation. Thus all the accumulated errors are less than
$\alpha$. This construction uses $n_A+n_B$ neurons, a polynomial number
in $K,M,\log(2m),\alpha^{-1}$.

\paragraph{Finite weights and genuine feedforward evaluation.}
Write the preceding covariance as
$Q'=\sum_{\ell=1}^r a_\ell h_\ell h_\ell^\top$, with
$a_\ell\le V:=\max\{1,A\}$ and $\|h_\ell\|_\infty\le B_h$.
A bulk preactivation is realized with weights
$\sqrt{a_\ell}\xi_\ell$, where the Gaussian vector $\xi$ is tilted by
a density ratio at most $e^{2GF_0}$. Intersecting the empirical covariance
concentration event with a Gaussian-coordinate tail event ensures
\[
 |\xi_\ell|\le R_\xi
 :=\sqrt{4GF_0+2\log(16n_B\max\{r,1\})}
\]
for every sampled coordinate. The latter event fails with probability
at most $1/8$, whereas the preceding covariance event fails with
probability less than $1/8$, so the intersection has positive probability.
Bulk weights are therefore bounded by $\sqrt V R_\xi$.
The atom weights are $a_\ell\sum_y\beta_j(y)h_\ell(y)$ and are bounded
by $VB_h/(4\eta)$. All new covariance coefficients are at most $V$:
they are $r_T(z)^2/n_B$ for bulk rows and $A/n_A$ for atom rows.

For example, if $B_W$ denotes the maximum of the two incoming-weight
bounds, the actual hidden coordinates satisfy the deterministic bound
$B_h^{\rm new}\le\lambda rB_WB_h$. This bounds the untruncated physical
activations as well as the weighted features. Hence at any fixed number
of layers both the weights and all physical hidden coordinates remain
polynomially bounded. The cutoff scalars, tilt matrices, and RKHS
certificates select finitely many real weights; they are not evaluated
on the input when the resulting network predicts.
At the first layer use input features $h_\ell(x)=\bar x_\ell$ and
coefficients $a_\ell=1/d$. In particular an atom row is the ordinary
linear form with weights $d^{-1}\sum_y\beta_j(y)\bar y$.

\subsection{A uniform functional exceptional-mass bound}
\begin{proposition}[Deep mean-field functional compression]
\label{prop:generic-meanfield-compression}
For each fixed $L$ and $\lambda$, there are polynomials
$P_{L,\lambda}$ and $J_{L,\lambda}$ such that the following holds for
every $d\ge1$, $K\ge1$, and
$N\ge P_{L,\lambda}(d,K)$. Except on an event of raw Gaussian probability
at most $(2H+1)e^{-NK}$, every globally strict-margin draw has the same
Boolean sign function as an auxiliary depth-$L$ network with normalized
energy at most $J_{L,\lambda}(d,K)$. All its neuronwise activations vanish
at zero and have Lipschitz constant at most $\lambda$.
\end{proposition}
\begin{proof}
Define layerwise deterministic diagonal bounds
\[
 M_s=2[8\lambda^2\{\log(4m)+A+1\}]^s,\qquad 0\le s\le H.
\]
At layer $s$, apply the one-step construction with $M=M_{s-1}$ and
its corresponding projection cutoff $B_s$. On its projection-good event,
the next covariance has diagonal at most
$2\lambda^2B_s^2+2\lambda^2M_{s-1}A\le M_s$. Starting from
$\max_xQ_0(x,x)\le2$, induction therefore validates the diagonal budget
used at each layer. Every $M_s$ is polynomial in $d,K$ for fixed
$L,\lambda$ because $\log m=d\log2$. There are at most two
exceptional events per hidden layer, as accounted for in the one-step
construction.

Set the desired final covariance error to $\zeta_H=1/(128K)$.
Working backwards, for $s=H,H-1,\ldots,1$, use the one-step construction
with $\alpha=\zeta_s$ and $M=M_{s-1}$, and define
$\zeta_{s-1}=r_s(\zeta_s)$ using that layer's modulus.
Each $\zeta_s^{-1}$ is polynomially bounded in $d,K$ at fixed $L$:
this is a fixed finite composition of the polynomial reciprocal moduli
above. Taking the maximum of the finitely many thresholds
$N_{0,s}(\zeta_s)$ defines a polynomial width threshold
$P_{L,\lambda}(d,K)$ (increased to at least one if necessary).
Starting from the exact $Q_0$, reconstruct the covariances forwards.
The transported certificates are fixed for the original covariances;
their estimates hold uniformly for every admissible preceding
reconstruction. Thus no independence assumption is imposed on those
certificates. Induction yields a weighted-feature covariance
$\widehat Q_H$ with
\[
 \|\widehat Q_H-Q_H\|_{\max}\le1/(128K).
\]
The construction above bounds every layer width, every physical weight,
and every physical activation by polynomials in $d,K$ at fixed depth.

For the actual Gaussian readout $v_i\sim N(0,N^{-2})$, put
$E_v=N\|v\|_2^2=N^{-1}\chi_N^2$. Markov's inequality at exponential
parameter $1/4$ gives
\[
 \Pr\{E_v>8K\}\le\exp[-N(2K-\tfrac12\log2)]\le e^{-NK}.
\]
On its complement the original output $z$ belongs to $\mathcal H_{Q_H}$
and $\|z\|_{Q_H}^2\le8K$. The scaled form of
Lemma~\ref{lem:rkhs-atom-certificate}, with approximation accuracy $1/8$,
gives
\[
 \|Q_H\beta-z\|_\infty\le1/8,\qquad \|\beta\|_1\le16K.
\]
Replacing $Q_H$ by $\widehat Q_H$ adds at most $1/8$ to this error.
For $\widehat Q_H=\sum_i a_i h_ih_i^\top$, the readout
$v_i'=2a_i\sum_y\beta(y)h_i(y)$ therefore represents a function with
the same signs as $z$ and strict margin greater than $1/2$.
Its weights are bounded by $32KVB_h$.

Pad the finitely many reconstructed layers to their maximum width $r_*$
with zero neurons, which is harmless since every activation vanishes at
zero. If $B_*$ bounds all weights, the normalized energy is at most
\[
 [d(d+1)+(L-1)r_*^2]B_*^2.
\]
Taking a polynomial upper bound on this expression defines
$J_{L,\lambda}(d,K)$. A conditional union bound over the $H$ layers and
the readout proves the stated raw exceptional probability. The auxiliary
capacity theorem applies uniformly to the possibly different activations;
no finite coding of activation formulas is needed.
\end{proof}
